\documentclass[UTF8,11pt]{ctexart}
\usepackage[a4paper,margin=24mm]{geometry}
\usepackage{amsmath,amssymb,amsthm,mathtools,mathrsfs}
\usepackage[all]{xy}
\usepackage{xurl,hyperref,enumitem}
\usepackage{tikz}
\usetikzlibrary{arrows.meta,cd,decorations.markings,calc,patterns}
\hypersetup{hidelinks,breaklinks=true}
\setlength{\emergencystretch}{3em}
\allowdisplaybreaks
\newtheorem*{theorem}{Theorem}
\newtheorem*{thm}{Theorem}
\newtheorem*{lemma}{Lemma}
\newtheorem*{proposition}{Proposition}
\newtheorem*{prop}{Proposition}
\newtheorem*{corollary}{Corollary}
\newtheorem*{cor}{Corollary}
\newtheorem*{claim}{Claim}
\theoremstyle{definition}
\newtheorem*{definition}{Definition}
\newtheorem*{defn}{Definition}
\newtheorem*{remark}{Remark}
\newtheorem*{example}{Example}
\newcommand{\R}{\mathbb R}
\newcommand{\C}{\mathbb C}
\newcommand{\Z}{\mathbb Z}
\newcommand{\Q}{\mathbb Q}
\newcommand{\N}{\mathbb N}
\newcommand{\F}{\mathbb F}
\newcommand{\D}{\mathbb D}
\newcommand{\bB}{\mathbb B}
\newcommand{\bC}{\mathbb C}
\newcommand{\bR}{\mathbb R}
\newcommand{\bZ}{\mathbb Z}
\newcommand{\bN}{\mathbb N}
\newcommand{\bQ}{\mathbb Q}
\newcommand{\bD}{\mathbb D}
\newcommand{\bF}{\mathbb F}
\newcommand{\CP}{\mathbb{CP}}
\newcommand{\RP}{\mathbb{RP}}
\newcommand{\sO}{\mathscr O}
\newcommand{\sI}{\mathscr I}
\newcommand{\sF}{\mathscr F}
\newcommand{\sG}{\mathscr G}
\newcommand{\sH}{\mathscr H}
\newcommand{\sR}{\mathscr R}
\newcommand{\sM}{\mathscr M}
\newcommand{\sV}{\mathscr V}
\newcommand{\sU}{\mathscr U}
\DeclareMathOperator{\Hom}{Hom}
\DeclareMathOperator{\Ext}{Ext}
\DeclareMathOperator{\Tor}{Tor}
\DeclareMathOperator{\Spec}{Spec}
\DeclareMathOperator{\Proj}{Proj}
\DeclareMathOperator{\supp}{supp}
\DeclareMathOperator{\im}{im}
\DeclareMathOperator{\id}{id}
\DeclareMathOperator{\Tr}{Tr}
\DeclareMathOperator{\tr}{tr}
\DeclareMathOperator{\rank}{rank}
\DeclareMathOperator{\ord}{ord}
\DeclareMathOperator{\dist}{dist}
\DeclareMathOperator{\End}{End}
\DeclareMathOperator{\Aut}{Aut}
\DeclareMathOperator{\Gal}{Gal}
\DeclareMathOperator{\vol}{vol}
\DeclareMathOperator{\Cl}{Cl}
\DeclareMathOperator{\coker}{coker}
\DeclareMathOperator{\codim}{codim}
\DeclareMathOperator{\divergence}{div}
\DeclareMathOperator{\Ric}{Ric}
\DeclareMathOperator{\grad}{grad}
\DeclareMathOperator{\Hess}{Hess}
\newcommand{\abs}[1]{\left\lvert#1\right\rvert}
\newcommand{\sabs}[1]{\lvert#1\rvert}
\newcommand{\babs}[1]{\bigl\lvert#1\bigr\rvert}
\newcommand{\Babs}[1]{\Bigl\lvert#1\Bigr\rvert}
\newcommand{\bbabs}[1]{\biggl\lvert#1\biggr\rvert}
\newcommand{\BBabs}[1]{\Biggl\lvert#1\Biggr\rvert}
\newcommand{\norm}[1]{\left\lVert#1\right\rVert}
\newcommand{\snorm}[1]{\lVert#1\rVert}
\newcommand{\bnorm}[1]{\bigl\lVert#1\bigr\rVert}
\newcommand{\Bnorm}[1]{\Bigl\lVert#1\Bigr\rVert}
\newcommand{\bbnorm}[1]{\biggl\lVert#1\biggr\rVert}
\newcommand{\BBnorm}[1]{\Biggl\lVert#1\Biggr\rVert}
\newcommand{\widebar}[1]{\overline{#1}}
\newcommand{\nicefrac}[2]{\frac{#1}{#2}}
\newcommand{\linnprod}[2]{\langle#1,#2\rangle}
\newcommand{\myindex}[1]{#1}
\newcommand{\glsadd}[1]{}
\newcommand{\avoidbreak}{}
\newcommand{\sourceref}[1]{\textnormal{[原文：\nolinkurl{#1}]}}
\newcommand{\thmref}[1]{\sourceref{#1}}
\newcommand{\lemmaref}[1]{\sourceref{#1}}
\newcommand{\propref}[1]{\sourceref{#1}}
\newcommand{\corref}[1]{\sourceref{#1}}
\newcommand{\defnref}[1]{\sourceref{#1}}
\newcommand{\exampleref}[1]{\sourceref{#1}}
\newcommand{\exerciseref}[1]{\sourceref{#1}}
\newcommand{\figureref}[1]{\sourceref{#1}}
\newcommand{\sectionref}[1]{\sourceref{#1}}
\newcommand{\appendixref}[1]{\sourceref{#1}}
\newcommand{\stacksref}[2]{\href{https://stacks.math.columbia.edu/tag/#1}{#2 [#1]}}


\newcommand{\E}{\mathbb E}
\newcommand{\Prob}{\mathbb P}
\newcommand{\Reff}{\mathcal R_{\mathrm{eff}}}
\newcommand{\eps}{\varepsilon}
\DeclareMathOperator{\diam}{diam}
\DeclareMathOperator{\Var}{Var}
\DeclareMathOperator{\Crit}{Crit}
\newcommand{\dd}{\,\mathrm d}


\begin{document}
\section*{089\quad 保测动力系统的Birkhoff逐点遍历收敛}
\subsection*{来源与整理范围}
Terence Tao，254A, Lecture 9: Ergodicity，2008-02-04；Theorem 1 的移位算子情形、Corollary 1 与 Theorem 2 的完整证明。
\par\url{https://terrytao.wordpress.com/2008/02/04/254a-lecture-9-ergodicity/}
\par\textbf{恢复方式（整理者说明）：}依据人类原文重新整理以补齐缺失题号；并非旧题证文件的逐字节恢复；2026-09-28。
\subsection*{作者命题与人类证明}

\paragraph{Notation.}
Let $(X,\mathcal X,\mu,T)$ be a measure-preserving probability system.
Write $Uf=f\circ T$ for the Koopman operator and
\[
 A_Nf=\frac1N\sum_{n=0}^{N-1}U^nf,
 \qquad \mathcal I=\{E\in\mathcal X:T^{-1}E=E\text{ modulo null sets}\}.
\]
\begin{theorem}[Pointwise ergodic theorem; Tao's Theorem 2]
For every $f\in L^1(X,\mathcal X,\mu)$,
\[
 A_Nf(x)\longrightarrow\E(f\mid\mathcal I)(x)
 \quad\text{for almost every }x.
\]
\end{theorem}

\begin{lemma}[Maximal inequality for the shift; Theorem 1 specialized to $U$]
For real $f\in L^1$, set $Mf=\sup_{N\geq1}A_Nf$. Then
\[
 \lambda\mu\{Mf>\lambda\}\leq\int_{\{Mf>\lambda\}}f\,d\mu
 \qquad(\lambda\in\R).
\]
Consequently, for arbitrary real or complex $f\in L^1$ and $\lambda>0$,
\[
 \mu\left\{\sup_N A_N|f|>\lambda\right\}\leq\frac{\|f\|_1}{\lambda}.
\]
\end{lemma}
\begin{proof}
Since $M(f-\lambda)=Mf-\lambda$, replace $f$ by $f-\lambda$ and reduce to
$\int_{\{Mf>0\}}f\,d\mu\geq0$. For $m\geq1$, define
\[
 F_m=\max_{0\leq N\leq m}\sum_{n=0}^{N-1}U^nf,
\]
where the sum for $N=0$ is zero. Then $F_m\geq0$, $F_m\in L^1$, and
\[
 F_m\leq F_{m+1}=\max(0,f+UF_m).
\]
On $\{F_m>0\}$ the second maximum equals $f+UF_m$, so integration and
positivity of $UF_m$ imply
\[
 \int_XF_m\,d\mu
 \leq\int_{\{F_m>0\}}f\,d\mu+\int_XUF_m\,d\mu.
\]
Measure preservation gives $\int UF_m=\int F_m$, and hence
$\int_{\{F_m>0\}}f\,d\mu\geq0$.
The sets $\{F_m>0\}$ increase to $\{Mf>0\}$. Dominated convergence,
with dominating function $|f|$, proves the first inequality. Apply it to
$|f|$ and use $\int_E|f|\leq\|f\|_1$ to obtain the stated maximal bound.
\end{proof}

\begin{proof}[Proof of the pointwise theorem]
Subtract $\E(f\mid\mathcal I)$, which is invariant, and reduce to proving
\begin{equation}\label{zero-pointwise}
 \limsup_{N\to\infty}|A_Nf(x)|=0\quad\text{almost everywhere}
\end{equation}
when $\E(f\mid\mathcal I)=0$. For a bounded coboundary $f=Ug-g$,
$g\in L^\infty$, the averages telescope:
\[
 A_Nf=\frac{U^Ng-g}{N}.
\]
This proves \eqref{zero-pointwise} for these functions.

The Hilbert-space argument used in the mean ergodic theorem shows that
bounded coboundaries are dense in the subspace of $L^2$ with zero
conditional expectation onto $\mathcal I$. That subspace is itself dense,
in $L^1$ norm, in the subspace of $L^1$ with zero conditional expectation.
Thus the pointwise assertion has been established for an $L^1$-dense class
in the latter subspace.

Choose $f_j$ in this dense class with $\|f-f_j\|_1\to0$. Outside a common
null set, $A_Nf_j(x)\to0$ for all $j$. The triangle inequality gives
\[
 \limsup_{N\to\infty}|A_Nf(x)|
 \leq\sup_N A_N|f-f_j|(x).
\]
By the maximal inequality, for any $\lambda>0$,
\[
 \mu\left\{\limsup_{N\to\infty}|A_Nf|>\lambda\right\}
 \leq\frac{\|f-f_j\|_1}{\lambda}.
\]
The left side does not depend on $j$; taking $j\to\infty$ makes it zero.
Taking positive rational $\lambda$ proves \eqref{zero-pointwise} and hence
the theorem.
\end{proof}

\subsection*{整理说明与先修范围（非作者正文）}
保留作者“有界余边界稠密类—极大不等式—$L^1$ 极限”的证明路线。原文 $T$ 同时表示点变换和函数算子，此处将后者记为 $U$。极大引理只取本题所用的移位算子，递推等式在该情形成立；不把该等式错误推广到一般正算子。原文 $L^1/L^2$ 稠密性说明中的 former/latter 指向已按相应范数关系纠正。

先修包括条件期望、均值遍历定理的 Hilbert 空间正交分解，以及有界函数的稠密性。正文引用的余边界稠密性来自均值遍历证明，不把逐点遍历定理本身当先修。本题保留完整极大不等式证明，从范数近似提升到几乎处处的轨道平均控制；未把 $L^1$ 范数收敛或其他更强遍历定理加入题面。
\subsection*{可修改的简短标注（非作者正文）}
$c$：保测动力系统中轨道平均的几乎处处极限

$a$：Koopman算子；条件期望；有界余边界；Hilbert空间正交分解；极大不等式；依测度控制

\texttt{cross\_theory\_status = yes}。将轨道平均转为函数空间上的算子平均，先借余边界望远镜求和处理稠密类，再用极大不等式把L1近似转回逐点结论；见两个证明的衔接。

全部标注为非穷尽、可修改初稿。
\subsection*{来源许可与编辑归属}
作者 About 页 Copyright etc. 允许带来源引用复用合理篇幅（例如单篇文章）；本题为原文指定部分的编辑转录。许可入口：https://terrytao.wordpress.com/about/ 。
ChatGPT 加入题号、中文来源说明、整理范围和初稿标注；编辑说明与人类证明分列。
\end{document}
